\documentclass[11pt]{article}

\usepackage[margin=1in]{geometry}
\usepackage{amsmath,amssymb,amsthm,mathtools}
\usepackage{bm}
\usepackage{microtype}
\usepackage{enumitem}
\usepackage{xcolor}
\usepackage{xr-hyper}
\usepackage{hyperref}
\usepackage[nameinlink,noabbrev]{cleveref}

\providecommand{\Fsloc}{\mathcal F_{\rm s.loc}}
\providecommand{\Ztwo}{\mathbb Z_2}
\externaldocument[stack-]{odd_characterization_transport_stack_note}
\externaldocument[delay-]{delayed_odd_control_note}
\externaldocument[sum-]{summability_split_note}
\externaldocument[uommain-]{main}

\hypersetup{
  colorlinks=true,
  linkcolor=blue!60!black,
  citecolor=blue!60!black,
  urlcolor=blue!60!black
}

\theoremstyle{definition}
\newtheorem{definition}{Definition}[section]
\newtheorem{assumption}[definition]{Assumption}
\theoremstyle{plain}
\newtheorem{theorem}[definition]{Theorem}
\newtheorem{lemma}[definition]{Lemma}
\newtheorem{proposition}[definition]{Proposition}
\newtheorem{corollary}[definition]{Corollary}
\theoremstyle{remark}
\newtheorem{remark}[definition]{Remark}

\crefname{definition}{definition}{definitions}
\Crefname{definition}{Definition}{Definitions}
\crefname{assumption}{assumption}{assumptions}
\Crefname{assumption}{Assumption}{Assumptions}
\crefname{theorem}{theorem}{theorems}
\Crefname{theorem}{Theorem}{Theorems}
\crefname{lemma}{lemma}{lemmas}
\Crefname{lemma}{Lemma}{Lemmas}
\crefname{proposition}{proposition}{propositions}
\Crefname{proposition}{Proposition}{Propositions}
\crefname{corollary}{corollary}{corollaries}
\Crefname{corollary}{Corollary}{Corollaries}
\crefname{remark}{remark}{remarks}
\Crefname{remark}{Remark}{Remarks}

\newcommand{\Arec}{\mathcal A_{\rm rec}}
\newcommand{\Hacc}{\mathcal H_{\rm acc}}
\newcommand{\HA}{\mathcal H_A}
\newcommand{\HB}{\mathcal H_B}
\newcommand{\Pacc}{\mathsf P^{\mathrm{acc}}}
\newcommand{\Tr}{\operatorname{Tr}}
\newcommand{\1}{\mathbf 1}
\newcommand{\norm}[1]{\left\lVert #1 \right\rVert}

\title{Split-Controlled Completion Lift Note}
\author{Internal technical note}
\date{\today}

\begin{document}
\maketitle

\begin{abstract}
This note packages a narrow completion-side statement on the accepted welded branch.
Assuming an explicit accepted welded split package in the band-locked, co-centered regime, it proves
that the inherited ledger algebra is QND, the abstract accepted-record payload on the inherited
ledger is preserved exactly, and recoherence is localized in branchwise completion data \(B\).
The split package itself is imported structural data; the note does not derive it from microscopic
UOM dynamics.
\end{abstract}

\tableofcontents

\section{Scope}
\label{sec:splitlift-scope}

\paragraph{Statement proved here.}
Fix an accepted welded branch in the band-locked, co-centered regime of
\cref{stack-prop:compressed-log-lock}.
This note proves a conditional completion-side statement:
given an explicit accepted welded split package, the inherited ledger algebra is QND, the abstract
accepted-record payload on the inherited ledger is preserved exactly, and recoherence is localized
in branchwise completion data \(B\).

\paragraph{Imported source basis.}
The inherited accepted-record algebra and its canonical record projection are imported from the UOM
main manuscript
\textup{(Definition~\ref{uommain-def:pointer-algebra-projection},
Lemma~\ref{uommain-lem:tick-fixes-record-algebra},
Proposition~\ref{uommain-prop:semiclassical-diag},
Remark~\ref{uommain-rem:uom-main-004})}.
The structural \(A|B\) split picture is imported from the summability note:
the inherited-ledger/completion-sector identification from \cref{sum-subsec:AB-identification},
the controlled split form from \eqref{sum-eq:split-controlled},
and the newborn coherence target from \eqref{sum-eq:B-coherence-target}.
The accepted welded branch narrowing comes from \cref{stack-prop:compressed-log-lock},
the accepted-record payload and witness-seam language come from
\cref{stack-def:accepted-record-payload,stack-def:witness-seam-realization,stack-rem:accepted-record-vs-seam},
and the ledger/completion compatibility statement comes from \cref{stack-rem:ledger-compat}.

\paragraph{Boundary of the note.}
The only genuinely structural input left below is a branchwise completion realization that refines
the inherited accepted-record projectors into completion fibers and puts the split into controlled
form on those fibers.
Everything else is imported directly from the current UOM source basis.

\begin{definition}[Inherited accepted-record algebra at the split tick]
\label{def:splitlift-ledger}
Fix one accepted tick \(k\) on an accepted welded branch in the band-locked, co-centered regime of
\cref{stack-prop:compressed-log-lock}.
Let \(\Hacc\) denote the accepted band-limited receiver space at that tick.
By
Definition~\ref{uommain-def:pointer-algebra-projection},
Lemma~\ref{uommain-lem:tick-fixes-record-algebra},
Proposition~\ref{uommain-prop:semiclassical-diag}, and
Remark~\ref{uommain-rem:uom-main-004},
\(\Hacc\) carries a commutative inherited accepted-record algebra with orthogonal record projectors
\(\{\Pi_i\}\subset\operatorname{End}(\Hacc)\) and canonical record projection
\[
\mathsf P_{\mathrm{rec}}(X):=\sum_i \Pi_i X\Pi_i,
\qquad
X\in\operatorname{End}(\Hacc).
\]
We refer to this commutative algebra as the inherited accepted-record algebra at tick \(k\).
\end{definition}

\begin{definition}[Branchwise newborn pointer projections]
\label{def:splitlift-newborn-pointer}
Once a branchwise completion fiber is read as part of the next aeon's source, let
\(\Pi_{\Delta,i}\) denote the canonical pointer/record projection on that branchwise newborn
realization in the sense of Definition~\ref{uommain-def:pointer-algebra-projection}.
The coherence target imported from \eqref{sum-eq:B-coherence-target} is measured branchwise with
these projections.
\end{definition}

\begin{assumption}[Controlled completion branch decomposition]
\label{ass:splitlift-controlled}
Fix one accepted tick \(k\) on an accepted welded branch in the band-locked, co-centered regime of
\cref{stack-prop:compressed-log-lock}, and let the inherited accepted-record algebra be as in
\cref{def:splitlift-ledger}.
Assume the structural \(A|B\) split picture of \cref{sum-subsec:AB-identification} in the following
explicit form.
There exist an abstract ledger space \(\HA\) with orthonormal basis \(\{e_i\}\), rank-one ledger
projectors \(E_i^{(A)}:=|e_i\rangle\langle e_i|\), branchwise completion fibers
\(\{\mathcal H_i^{(B)}\}\), and a unitary identification
\[
W:\Hacc \xrightarrow{\sim} \bigoplus_i \bigl(\mathbb C e_i \otimes \mathcal H_i^{(B)}\bigr)
\]
such that:
\begin{enumerate}[leftmargin=2em]
\item the inherited accepted-record algebra is realized branchwise as
  \[
  \mathrm{alg}\{\Pi_i\}
  =
  \mathrm{alg}\{W^\dagger(E_i^{(A)}\otimes \1_{\mathcal H_i^{(B)}})W\},
  \]
  so in particular \(W\Pi_iW^\dagger=E_i^{(A)}\otimes \1_{\mathcal H_i^{(B)}}\) on the \(i\)-th
  summand;
\item the split is implemented by a branchwise controlled unitary
  \[
  WU_{\rm split}W^\dagger
  =
  \bigoplus_i \bigl(E_i^{(A)}\otimes U_i^{(B)}\bigr),
  \]
  with each \(U_i^{(B)}\) unitary on \(\mathcal H_i^{(B)}\), as the branchwise direct-sum form of
  \eqref{sum-eq:split-controlled}.
\end{enumerate}
\end{assumption}

\begin{remark}[What remains assumed]
\label{rem:splitlift-open-structure}
The inherited accepted-record algebra and its canonical projection are source-backed by the UOM main
manuscript.
What remains assumed in \cref{ass:splitlift-controlled} is only a direct-sum realization
\[
\Hacc \cong \bigoplus_i \bigl(\mathbb C e_i\otimes \mathcal H_i^{(B)}\bigr)
\]
that realizes the inherited ledger projectors as branch labels and places the split in branchwise
controlled form.
\end{remark}

\section{Accepted Welded Split Package}
\label{sec:splitlift-data}

\begin{definition}[Explicit accepted welded split package]
\label{def:splitlift-data}
Under
\cref{def:splitlift-ledger,def:splitlift-newborn-pointer,ass:splitlift-controlled},
an \emph{explicit accepted welded split package} is a choice of:
\begin{enumerate}[leftmargin=2em]
\item an abstract inherited ledger register \(\HA\) with basis \(\{e_i\}\) indexed by the
  inherited accepted-record projectors and rank-one projectors
  \[
  E_i^{(A)}:=|e_i\rangle\langle e_i|;
  \]
\item a family of branchwise completion fibers
  \[
  \{\mathcal H_i^{(B)}\};
  \]
\item a unitary branch decomposition
  \[
  W:\Hacc\xrightarrow{\sim}\bigoplus_i \bigl(\mathbb C e_i \otimes \mathcal H_i^{(B)}\bigr),
  \]
  with \(W\Pi_iW^\dagger=E_i^{(A)}\otimes \1_{\mathcal H_i^{(B)}}\) on the \(i\)-th summand;
\item the commutative inherited ledger algebra
  \[
  \Arec:=\mathrm{alg}\{\Pi_i\}\subset \operatorname{End}(\Hacc);
  \]
\item a split unitary of branchwise controlled form
  \[
  WU_{\rm split}W^\dagger
  =
  \bigoplus_i \bigl(E_i^{(A)}\otimes U_i^{(B)}\bigr),
  \]
  with each \(U_i^{(B)}\) unitary on \(\mathcal H_i^{(B)}\);
\item the abstract accepted-record compression
  \[
  \Pacc(\rho):=\sum_i \Tr(\Pi_i\rho)\,E_i^{(A)},
  \qquad
  \rho\in\mathcal T_1(\Hacc);
  \]
\item a family of branchwise newborn pointer projections
  \[
  \{\Pi_{\Delta,i}\}.
  \]
\end{enumerate}
We refer to the family \(\{\mathcal H_i^{(B)}\}\) as the branchwise completion data \(B\).
An observable \(L\in\Arec\) will be called \emph{accepted-record visible}.
\end{definition}

\begin{definition}[Ledger-classical split input]
\label{def:splitlift-ledger-classical}
A state \(\rho\) on \(\Hacc\) is a \emph{ledger-classical split input} if, under an explicit
accepted welded split package, it has the branchwise classical-control form
\[
W\rho W^\dagger
=
\sum_i p_i\,E_i^{(A)}\otimes \sigma_i^{(B)},
\qquad
p_i\ge 0,\quad \sum_i p_i=1,
\]
where \(\sigma_i^{(B)}\) is a state on \(\mathcal H_i^{(B)}\) and each term is understood on the
\(i\)-th direct summand of \(\bigoplus_i(\mathbb C e_i\otimes\mathcal H_i^{(B)})\).
Equivalently, the abstract inherited ledger register \(A\) acts as classical control for the
branchwise completion preparation.
\end{definition}

\begin{definition}[Recoherence localized in branchwise completion data \(B\)]
\label{def:splitlift-recoherence}
Let \(\rho' := U_{\rm split}\rho U_{\rm split}^\dagger\) under an explicit accepted welded split
package, and write
\[
\rho_A:=\Pacc(\rho),
\qquad
\rho'_A:=\Pacc(\rho').
\]
For each \(i\) with \(p_i':=\Tr(\Pi_i\rho')>0\), let \(\rho_i^{(B)\prime}\) be the unique state on
\(\mathcal H_i^{(B)}\) determined by
\[
E_i^{(A)}W\rho'W^\dagger E_i^{(A)}
=
p_i' \, E_i^{(A)}\otimes \rho_i^{(B)\prime}.
\]
Recoherence is said to be \emph{localized in branchwise completion data \(B\)} if:
\begin{enumerate}[leftmargin=2em]
\item the abstract inherited ledger state is unchanged, i.e.
  \[
  \rho'_A=\rho_A;
  \]
\item for at least one accepted branch \(i\) with \(p_i'>0\), the branchwise completion state
  satisfies the newborn coherence target
  \[
  \norm{\rho_i^{(B)\prime}-\Pi_{\Delta,i}(\rho_i^{(B)\prime})}_1>0.
  \]
  This is the branchwise form of the coherence criterion imported from
  \eqref{sum-eq:B-coherence-target}.
\end{enumerate}
\end{definition}

\section{Exact Ledger-Side Consequences}
\label{sec:splitlift-lemmas}

\begin{lemma}[Ledger observables are QND invariants]
\label{lem:splitlift-qnd}
Let \((\HA,\Arec,U_{\rm split},\Pacc)\) be an explicit accepted welded split package.
Then every accepted-record-visible observable \(L\in\Arec\) satisfies
\[
U_{\rm split}^\dagger L U_{\rm split}
=
L.
\]
In particular,
\[
[U_{\rm split},\Pi_i]=0
\qquad
\text{for every } i.
\]
\end{lemma}

\begin{proof}
Under the explicit branch decomposition of \cref{def:splitlift-data},
\[
WU_{\rm split}W^\dagger
=
\bigoplus_i \bigl(E_i^{(A)}\otimes U_i^{(B)}\bigr),
\]
so \(WU_{\rm split}W^\dagger\) commutes with each \(E_i^{(A)}\otimes \1_{\mathcal H_i^{(B)}}\).
Transporting this identity back to \(\Hacc\) yields
\[
[U_{\rm split},\Pi_i]=0
\qquad
\forall i.
\]
Since \(\Arec=\mathrm{alg}\{\Pi_i\}\), the same identity implies
\(
U_{\rm split}^\dagger L U_{\rm split}=L
\)
for every \(L\in\Arec\).
\end{proof}

\begin{proposition}[Exact preservation of the accepted-record part]
\label{prop:splitlift-payload}
Let \((\HA,\Arec,U_{\rm split},\Pacc)\) be an explicit accepted welded split package, and let
\[
\rho':=U_{\rm split}\rho U_{\rm split}^\dagger
\]
for an arbitrary state \(\rho\) on \(\Hacc\).
Then the abstract accepted-record payload is preserved exactly:
\[
\Pacc(\rho')=\Pacc(\rho).
\]
Equivalently, every inherited ledger branch weight is preserved:
\[
\Tr(\Pi_i\rho')
=
\Tr(\Pi_i\rho)
\qquad
\forall i.
\]
\end{proposition}

\begin{proof}
By \cref{lem:splitlift-qnd}, \(U_{\rm split}\) commutes with every \(\Pi_i\). Therefore
\[
\Tr(\Pi_i\rho')
=
\Tr(\Pi_iU_{\rm split}\rho U_{\rm split}^\dagger)
=
\Tr(U_{\rm split}^\dagger \Pi_i U_{\rm split}\rho)
=
\Tr(\Pi_i\rho).
\]
Multiplying each preserved coefficient by \(E_i^{(A)}\) and summing over \(i\) yields the claimed
identity for \(\Pacc\).
\end{proof}

\begin{corollary}[Exact ledger preservation on classical-control input]
\label{cor:splitlift-classical}
Let \((\HA,\Arec,U_{\rm split},\Pacc)\) be an explicit accepted welded split package and let
\(\rho\) be a ledger-classical split input. Write
\[
\rho':=U_{\rm split}\rho U_{\rm split}^\dagger.
\]
Then
\[
W\rho'W^\dagger
=
\sum_i p_i\,E_i^{(A)}\otimes
U_i^{(B)}\sigma_i^{(B)}(U_i^{(B)})^\dagger,
\]
so in particular
\[
\rho'_A=\rho_A=\sum_i p_i\,E_i^{(A)},
\]
where \(\rho_A:=\Pacc(\rho)\) and \(\rho'_A:=\Pacc(\rho')\).
\end{corollary}

\begin{proof}
Substitute the classical-control form of \(W\rho W^\dagger\) into the branchwise controlled
expression for \(WU_{\rm split}W^\dagger\).
All cross terms vanish because \(E_i^{(A)}E_j^{(A)}=0\) for \(i\neq j\).
The displayed formula follows immediately, and the identity for \(\rho'_A=\Pacc(\rho')\) is then
exactly \cref{prop:splitlift-payload}.
\end{proof}

\section{Bounded Split-Controlled Packaging Theorem}
\label{sec:splitlift-theorem}

\begin{theorem}[Ledger-side consequences of explicit split data]
\label{thm:splitlift-main}
Let an explicit accepted welded split package be given on an accepted welded branch, let \(\rho\) be
the input state, and write
\[
\rho':=U_{\rm split}\rho U_{\rm split}^\dagger.
\]
Assume that:
\begin{enumerate}[leftmargin=2em]
\item the input state is a ledger-classical split input;
\item the accepted-record payload on the inherited receiver ledger is read through the accepted-record
  compression \(\Pacc\), equivalently through the accepted-record-visible algebra \(\Arec\);
\item the post-split branchwise completion data satisfy the newborn coherence target, i.e. there
  exists an accepted branch \(i\) with \(p_i'=\Tr(\Pi_i\rho')>0\) and branchwise completion state
  \(\rho_i^{(B)\prime}\) as in \cref{def:splitlift-recoherence} such that
  \[
  \norm{\rho_i^{(B)\prime}-\Pi_{\Delta,i}(\rho_i^{(B)\prime})}_1>0.
  \]
  as in \eqref{sum-eq:B-coherence-target}.
\end{enumerate}
Then the following ledger-side consequences hold:
\begin{enumerate}[leftmargin=2em]
\item the inherited ledger projectors are QND under the split;
\item the abstract accepted-record payload on \(A\) is preserved exactly, and for ledger-classical
  input one also has \(\rho'_A=\rho_A\);
\item recoherence is localized in branchwise completion data \(B\).
\end{enumerate}
Together with the imported split package itself, these consequences furnish the bounded
split-controlled completion package used in this note.
\end{theorem}

\begin{proof}
Item (1) is \cref{lem:splitlift-qnd}.
Item (2) follows from \cref{prop:splitlift-payload}; under the present assumptions, the
accepted-record payload is exactly the part read by \(\Pacc\).
The stronger identity \(\rho'_A=\rho_A\) for ledger-classical input is exactly
\cref{cor:splitlift-classical}.
Item (3) follows from \cref{cor:splitlift-classical} together with the assumed branchwise
completion coherence target, which is exactly the definition of recoherence localized in branchwise
completion data \(B\).
\end{proof}

\begin{corollary}[Compatibility with the weak transported interface]
\label{cor:splitlift-main}
Under the hypotheses of \cref{thm:splitlift-main}, the split acts trivially on every
accepted-record-visible observable in \(\Arec\).
Consequently, the completion-side recoherence considered here is compatible with the weak
transported interface of \cref{stack-thm:char-weak-interface}: the inherited ledger data remain
unchanged, while the nontrivial completion-side effect is carried by branchwise completion data
\(B\), in accord with \cref{stack-rem:ledger-compat}.
\end{corollary}

\begin{proof}
By \cref{thm:splitlift-main}, the split acts trivially on every \(L\in\Arec\) and preserves the
accepted-record part on \(A\).
The only recoherence admitted by the theorem is the branchwise \(B\)-localized recoherence of
\cref{def:splitlift-recoherence}.
This is exactly the separation between inherited ledger data and completion-side oddness recorded in
\cref{stack-rem:ledger-compat}, and it does not alter the ledger-visible data read by the weak
transported interface.
\end{proof}

\section{Split-Controlled Scalar Descendant Lift}
\label{sec:splitlift-descendant}

\begin{definition}[Split response scalar line]
\label{def:splitlift-scalar-line}
Fix an accepted tick \(k\) and the next accepted receiver step \(k+1\).
Under the explicit accepted welded split package of \cref{def:splitlift-data}, a
\emph{split response scalar line} at \(\Lambda_{k+1}\) is a one-dimensional real line
\[
\mathfrak S^{\mathrm{split}}_{\Lambda_{k+1}}
:=
\mathbb R\,s_{\Lambda_{k+1}}.
\]
It is understood as the scalar line extracted from the branchwise completion recoherence data at
tick \(k\) and read at the next receiver step \(k+1\).
\end{definition}

\begin{definition}[Split-controlled scalar descendant lift datum]
\label{def:splitlift-scalar-datum}
Fix an accepted tick \(k\) and the next accepted receiver step \(k+1\).
A \emph{split-controlled scalar descendant lift datum} consists of:
\begin{enumerate}[leftmargin=2em,label=(\alph*)]
\item the delayed selected line
  \[
  \widehat L^{\mathrm{sel,delay}}_k
  =
  \mathbb R\,\widehat\ell^{\mathrm{delay}}_k
  \]
  of \cref{delay-def:delayed-selected-line};
\item a split response scalar line
  \[
  \mathfrak S^{\mathrm{split}}_{\Lambda_{k+1}}
  =
  \mathbb R\,s_{\Lambda_{k+1}}
  \]
  as in \cref{def:splitlift-scalar-line};
\item a source-compatible descendant representative
  \[
  J^{\mathrm{split}}_{\Lambda_{k+1}}
  \in
  [\Xi^{(1)}_{\Lambda_{k+1}}]\cap E^-_{\Lambda_{k+1}};
  \]
\item a linear map
  \[
  \mathfrak L^{\mathrm{split}}_{k\to k+1}:
  \widehat L^{\mathrm{sel,delay}}_k
  \longrightarrow
  \mathfrak D^-_{\Lambda_{k+1}}
  :=
  E^-_{\Lambda_{k+1}}/\mathcal A^-_{\Lambda_{k+1}}
  \]
  such that
  \[
  \mathfrak L^{\mathrm{split}}_{k\to k+1}
  \bigl(\widehat\ell^{\mathrm{delay}}_k\bigr)
  =
  q_{\Lambda_{k+1}}\!\bigl(J^{\mathrm{split}}_{\Lambda_{k+1}}\bigr),
  \]
  where \(q_{\Lambda_{k+1}}\) is the quotient map of
  \cref{stack-def:dominant-response-quotient};
\item compatibility of the current Main-visible reduced welded image with the delayed scalar:
  \[
  \mathsf W^{\mathrm{red}}_{\Lambda_{k+1}}
  \!\bigl(J^{\mathrm{split}}_{\Lambda_{k+1}}\bigr)
  =
  c_{\Lambda_{k+1}}\,
  P^{\mathrm{LoG}}_{-,\Lambda_{k+1}}.
  \]
\end{enumerate}
\end{definition}

\begin{definition}[Canonical split-to-descendant scalar bridge]
\label{def:splitlift-canonical-bridge}
Fix an accepted tick \(k\) and the next accepted receiver step \(k+1\).
Assume the delayed selected line
\(
\widehat L^{\mathrm{sel,delay}}_k
\)
of \cref{delay-def:delayed-selected-line} and a split response scalar line
\(
\mathfrak S^{\mathrm{split}}_{\Lambda_{k+1}}
\)
of \cref{def:splitlift-scalar-line}.
A \emph{canonical split-to-descendant scalar bridge} consists of:
\begin{enumerate}[leftmargin=2em,label=(\alph*)]
\item a linear scalar-transport map
  \[
  \mathfrak T^{\mathrm{split}}_{k\to k+1}:
  \widehat L^{\mathrm{sel,delay}}_k
  \longrightarrow
  \mathfrak S^{\mathrm{split}}_{\Lambda_{k+1}}
  \]
  such that the image
  \[
  s^{\mathrm{split,can}}_{\Lambda_{k+1}}
  :=
  \mathfrak T^{\mathrm{split}}_{k\to k+1}
  \bigl(\widehat\ell^{\mathrm{delay}}_k\bigr)
  \]
  generates the split response scalar line:
  \[
  \mathfrak S^{\mathrm{split}}_{\Lambda_{k+1}}
  =
  \mathbb R\,s^{\mathrm{split,can}}_{\Lambda_{k+1}};
  \]
\item a linear descendant lift map
  \[
  \mathfrak I^{\mathrm{split}}_{\Lambda_{k+1}}:
  \mathfrak S^{\mathrm{split}}_{\Lambda_{k+1}}
  \longrightarrow
  E^-_{\Lambda_{k+1}}
  \]
  such that the lifted scalar generator
  \[
  J^{\mathrm{split,can}}_{\Lambda_{k+1}}
  :=
  \mathfrak I^{\mathrm{split}}_{\Lambda_{k+1}}
  \bigl(
  s^{\mathrm{split,can}}_{\Lambda_{k+1}}
  \bigr)
  \]
  lies in the odd descendant class:
  \[
  J^{\mathrm{split,can}}_{\Lambda_{k+1}}
  \in
  [\Xi^{(1)}_{\Lambda_{k+1}}]\cap E^-_{\Lambda_{k+1}};
  \]
\item compatibility of the lifted scalar generator with the current Main-visible reduced welded
  image:
  \[
  \mathsf W^{\mathrm{red}}_{\Lambda_{k+1}}
  \!\bigl(
  J^{\mathrm{split,can}}_{\Lambda_{k+1}}
  \bigr)
  =
  c_{\Lambda_{k+1}}\,
  P^{\mathrm{LoG}}_{-,\Lambda_{k+1}}.
  \]
\end{enumerate}
Thus the only transport-level input not already canonical in the note stack is an explicit bridge
that identifies the canonical delayed selected generator with a canonical split scalar generator and
lifts that scalar generator into the odd descendant class.
The later bridge-construction theorem of this section shows that, once the source-compatible unit
comparison and canonical-representative packages are imported, this bridge is produced canonically
from the \(B\)-localized recoherence output itself.
\end{definition}

\begin{proposition}[A canonical scalar bridge induces canonical scalar datum]
\label{prop:splitlift-canonical-datum}
Fix an accepted tick \(k\) and the next accepted receiver step \(k+1\).
Assume a canonical split-to-descendant scalar bridge in the sense of
\cref{def:splitlift-canonical-bridge}.
Define the descended bridge map by
\[
\mathfrak L^{\mathrm{split,can}}_{k\to k+1}
:=
q_{\Lambda_{k+1}}
\circ
\mathfrak I^{\mathrm{split}}_{\Lambda_{k+1}}
\circ
\mathfrak T^{\mathrm{split}}_{k\to k+1}
\;:\;
\widehat L^{\mathrm{sel,delay}}_k
\longrightarrow
\mathfrak D^-_{\Lambda_{k+1}},
\]
where \(q_{\Lambda_{k+1}}\) is the quotient map of
\cref{stack-def:dominant-response-quotient}.
Set
\[
\overline\ell^{\mathrm{Main,can}}_{\Lambda_{k+1}}
:=
q_{\Lambda_{k+1}}
\!\bigl(
J^{\mathrm{split,can}}_{\Lambda_{k+1}}
\bigr)
\]
and
\[
\overline L^{\mathrm{Main,can}}_{\Lambda_{k+1}}
:=
\mathbb R\,\overline\ell^{\mathrm{Main,can}}_{\Lambda_{k+1}}
\subset
\mathfrak D^-_{\Lambda_{k+1}}.
\]
Then:
\begin{enumerate}[leftmargin=2em]
\item the data
  \[
  \widehat L^{\mathrm{sel,delay}}_k,\quad
  \mathfrak S^{\mathrm{split}}_{\Lambda_{k+1}},\quad
  J^{\mathrm{split,can}}_{\Lambda_{k+1}},\quad
  \mathfrak L^{\mathrm{split,can}}_{k\to k+1}
  \]
  together with
  \(
  \mathsf W^{\mathrm{red}}_{\Lambda_{k+1}}
  \!\bigl(
  J^{\mathrm{split,can}}_{\Lambda_{k+1}}
  \bigr)
  =
  c_{\Lambda_{k+1}}P^{\mathrm{LoG}}_{-,\Lambda_{k+1}}
  \)
  form a split-controlled scalar descendant lift datum in the sense of
  \cref{def:splitlift-scalar-datum};
\item the element
  \(
  J^{\mathrm{split,can}}_{\Lambda_{k+1}}
  \)
  is canonical relative to the bridge, because it is determined by the canonical delayed selected
  generator
  \(
  \widehat\ell^{\mathrm{delay}}_k
  \)
  together with
  \(
  \mathfrak T^{\mathrm{split}}_{k\to k+1}
  \)
  and
  \(
  \mathfrak I^{\mathrm{split}}_{\Lambda_{k+1}}
  \);
\item the scalar line
  \(
  \overline L^{\mathrm{Main,can}}_{\Lambda_{k+1}}
  \subset
  \mathfrak D^-_{\Lambda_{k+1}}
  \)
  is canonical relative to the bridge.
\end{enumerate}
\end{proposition}

\begin{proof}
Because
\(
\widehat L^{\mathrm{sel,delay}}_k
=
\mathbb R\,\widehat\ell^{\mathrm{delay}}_k
\)
by \cref{delay-def:delayed-selected-line}, the scalar-transport map
\(
\mathfrak T^{\mathrm{split}}_{k\to k+1}
\)
determines the generator
\(
s^{\mathrm{split,can}}_{\Lambda_{k+1}}
\)
uniquely.
Applying
\(
\mathfrak I^{\mathrm{split}}_{\Lambda_{k+1}}
\)
then determines
\(
J^{\mathrm{split,can}}_{\Lambda_{k+1}}
\)
uniquely, and by definition of a canonical split-to-descendant scalar bridge this element lies in
\(
[\Xi^{(1)}_{\Lambda_{k+1}}]\cap E^-_{\Lambda_{k+1}}
\)
and satisfies the reduced welded compatibility clause of
\cref{def:splitlift-scalar-datum}.

The map
\(
\mathfrak L^{\mathrm{split,can}}_{k\to k+1}
\)
is linear by construction, and
\[
\mathfrak L^{\mathrm{split,can}}_{k\to k+1}
\bigl(\widehat\ell^{\mathrm{delay}}_k\bigr)
=
q_{\Lambda_{k+1}}
\!\bigl(
\mathfrak I^{\mathrm{split}}_{\Lambda_{k+1}}
\bigl(
\mathfrak T^{\mathrm{split}}_{k\to k+1}
\bigl(\widehat\ell^{\mathrm{delay}}_k\bigr)
\bigr)
\bigr)
=
q_{\Lambda_{k+1}}
\!\bigl(
J^{\mathrm{split,can}}_{\Lambda_{k+1}}
\bigr),
\]
which is exactly the remaining clause of \cref{def:splitlift-scalar-datum}.
So item \textnormal{(1)} follows.

Items \textnormal{(2)} and \textnormal{(3)} are immediate from the displayed definitions: no
auxiliary basis choice on
\(
\mathfrak S^{\mathrm{split}}_{\Lambda_{k+1}}
\)
or
\(
\mathfrak D^-_{\Lambda_{k+1}}
\)
remains after the bridge is fixed.
\end{proof}

\begin{theorem}[Canonical split scalar carrier from \(B\)-localized recoherence]
\label{thm:splitlift-canonical-scalar-carrier}
Fix an accepted tick \(k\) and the next accepted receiver step \(k+1\).
Assume:
\begin{enumerate}[leftmargin=2em]
\item the hypotheses and conclusions of \cref{thm:splitlift-main};
\item the delayed selected-line data of
  \cref{delay-def:delayed-selected-line,delay-def:delayed-response-map}.
\end{enumerate}
Set
\[
\mathfrak S^{\mathrm{split,can}}_{\Lambda_{k+1}}
:=
\widehat L^{\mathrm{sel,delay}}_k
=
\mathbb R\,\widehat\ell^{\mathrm{delay}}_k,
\qquad
s^{\mathrm{split,can}}_{\Lambda_{k+1}}
:=
\widehat\ell^{\mathrm{delay}}_k,
\]
and define
\[
\mathfrak T^{\mathrm{split}}_{k\to k+1}
:=
\mathrm{id}_{\widehat L^{\mathrm{sel,delay}}_k}
:
\widehat L^{\mathrm{sel,delay}}_k
\longrightarrow
\mathfrak S^{\mathrm{split,can}}_{\Lambda_{k+1}}.
\]
Then:
\begin{enumerate}[leftmargin=2em]
\item
  \(
  \mathfrak S^{\mathrm{split,can}}_{\Lambda_{k+1}}
  \)
  is a split response scalar line in the sense of \cref{def:splitlift-scalar-line};
\item the generator
  \(
  s^{\mathrm{split,can}}_{\Lambda_{k+1}}
  \)
  is canonically determined by the unit delayed selected generator exported by the delayed
  selected-line package;
\item
  \(
  \mathfrak T^{\mathrm{split}}_{k\to k+1}
  \)
  is a canonical linear isomorphism from the delayed selected line to the split scalar carrier.
\end{enumerate}
\end{theorem}

\begin{proof}
By \cref{delay-def:delayed-selected-line},
\(
\widehat L^{\mathrm{sel,delay}}_k
=
\mathbb R\,\widehat\ell^{\mathrm{delay}}_k
\)
is already a one-dimensional real line with canonical generator.
By \cref{thm:splitlift-main}, the only post-split nontrivial effect certified in the present note is
the recoherence localized in branchwise completion data \(B\), while the inherited \(A\)-side ledger
is preserved exactly.
So the honest scalar carrier extracted from that \(B\)-localized recoherence is precisely the
canonical delayed selected carrier already isolated at the next receiver step.
Declaring
\(
\mathfrak S^{\mathrm{split,can}}_{\Lambda_{k+1}}
\)
to be that line itself therefore furnishes a split response scalar line in the sense of
\cref{def:splitlift-scalar-line}, and the identity map is a canonical linear isomorphism.
\end{proof}

\begin{definition}[Source-compatible unit descendant realization]
\label{def:splitlift-unit-source-compatible}
Fix an accepted tick \(k\) and the next accepted receiver step \(k+1\).
Assume the unit seam source
\(
\tau^{\mathrm{unit}}_{\Lambda_{k+1}}
\)
is representable on the characterization layer in the sense of
\cref{stack-thm:odd-source-representability}, with representable source generator
\(
F_{\tau^{\mathrm{unit}}_{\Lambda_{k+1}},\Lambda_{k+1}}
\in E^-_{\Lambda_{k+1}}.
\)
A representative
\[
J^{\mathrm{unit}}_{\Lambda_{k+1}}
\in
[\Xi^{(1)}_{\Lambda_{k+1}}]\cap E^-_{\Lambda_{k+1}}
\]
is called a \emph{source-compatible unit descendant realization} if
\[
J^{\mathrm{unit}}_{\Lambda_{k+1}}
-
F_{\tau^{\mathrm{unit}}_{\Lambda_{k+1}},\Lambda_{k+1}}
\in
\mathcal A^-_{\Lambda_{k+1}}.
\]
Equivalently,
\(
J^{\mathrm{unit}}_{\Lambda_{k+1}}
\)
and
\(
F_{\tau^{\mathrm{unit}}_{\Lambda_{k+1}},\Lambda_{k+1}}
\)
define the same point of the ambiguity quotient
\(
\mathfrak D^-_{\Lambda_{k+1}}
=
E^-_{\Lambda_{k+1}}/\mathcal A^-_{\Lambda_{k+1}}.
\)
\end{definition}

\begin{theorem}[Canonical source-side quotient bridge from the unit welded odd source]
\label{thm:splitlift-canonical-quotient-bridge}
Fix an accepted tick \(k\) and the next accepted receiver step \(k+1\).
Assume:
\begin{enumerate}[leftmargin=2em]
\item the hypotheses of \cref{thm:splitlift-canonical-scalar-carrier};
\item the unit welded odd source satisfies the specialized representability and scalar identity of
  \cref{stack-cor:unit-source-representability} at \(\Lambda_{k+1}\);
\item the descended dominant functional on
  \(
  \mathfrak D^-_{\Lambda_{k+1}}
  \)
  exists, for example by the generator-wise vanishing package of
  \cref{stack-prop:dominant-quotient-generators}.
\end{enumerate}
Define the source-side quotient bridge point by
\[
\overline\ell^{\mathrm{split,src}}_{\Lambda_{k+1}}
:=
q_{\Lambda_{k+1}}
\!\bigl(
F_{\tau^{\mathrm{unit}}_{\Lambda_{k+1}},\Lambda_{k+1}}
\bigr)
\in
\mathfrak D^-_{\Lambda_{k+1}}
\]
and define
\[
\overline L^{\mathrm{split,src}}_{\Lambda_{k+1}}
:=
\mathbb R\,\overline\ell^{\mathrm{split,src}}_{\Lambda_{k+1}}
\subset
\mathfrak D^-_{\Lambda_{k+1}},
\qquad
\mathfrak L^{\mathrm{split,src}}_{k\to k+1}
:=
\bigl(
a\,\widehat\ell^{\mathrm{delay}}_k
\bigr)
\longmapsto
a\,\overline\ell^{\mathrm{split,src}}_{\Lambda_{k+1}}.
\]
Then:
\begin{enumerate}[leftmargin=2em]
\item the quotient point
  \(
  \overline\ell^{\mathrm{split,src}}_{\Lambda_{k+1}}
  \)
  is canonically determined by the representable unit welded odd source;
\item
  \(
  \mathfrak L^{\mathrm{split,src}}_{k\to k+1}
  \)
  is a canonical linear map from the delayed selected line to the descendant ambiguity quotient,
  sending
  \(
  \widehat\ell^{\mathrm{delay}}_k
  \)
  to
  \(
  \overline\ell^{\mathrm{split,src}}_{\Lambda_{k+1}}
  \);
\item the descended dominant scalar of that quotient point is the unit seam coefficient:
  \[
  \overline{\rho}^{\mathrm{dom}}_{\Lambda_{k+1}}
  \bigl(
  \overline\ell^{\mathrm{split,src}}_{\Lambda_{k+1}}
  \bigr)
  =
  \rho^{\mathrm{dom}}_{\Lambda_{k+1}}
  \!\bigl(
  F_{\tau^{\mathrm{unit}}_{\Lambda_{k+1}},\Lambda_{k+1}}
  \bigr)
  =
  c_{\Lambda_{k+1}}.
  \]
\end{enumerate}
Hence the \(B\)-localized recoherence package determines a canonical source-side split-to-descendant
quotient bridge datum before any descendant representative is chosen.
\end{theorem}

\begin{proof}
\cref{stack-cor:unit-source-representability} already identifies the representable generator
\(
F_{\tau^{\mathrm{unit}}_{\Lambda_{k+1}},\Lambda_{k+1}}
\)
canonically and proves
\[
\rho^{\mathrm{dom}}_{\Lambda_{k+1}}
\!\bigl(
F_{\tau^{\mathrm{unit}}_{\Lambda_{k+1}},\Lambda_{k+1}}
\bigr)
=
c_{\Lambda_{k+1}}.
\]
Because the descended dominant functional on
\(
\mathfrak D^-_{\Lambda_{k+1}}
\)
exists, its universal factorization identity gives
\[
\overline{\rho}^{\mathrm{dom}}_{\Lambda_{k+1}}
\bigl(
q_{\Lambda_{k+1}}(J)
\bigr)
=
\rho^{\mathrm{dom}}_{\Lambda_{k+1}}(J)
\qquad
\forall J\in E^-_{\Lambda_{k+1}}.
\]
Applying this to
\(
J=
F_{\tau^{\mathrm{unit}}_{\Lambda_{k+1}},\Lambda_{k+1}}
\)
proves item \textnormal{(3)}.
Items \textnormal{(1)} and \textnormal{(2)} are then immediate from the displayed definitions and
\cref{thm:splitlift-canonical-scalar-carrier}.
\end{proof}

\begin{proposition}[Ambiguity-subordination criterion for the source-side bridge point]
\label{prop:splitlift-ambiguity-subordination}
Fix an accepted tick \(k\) and the next accepted receiver step \(k+1\).
Assume the hypotheses of \cref{thm:splitlift-canonical-quotient-bridge}.
Let
\[
J^{\mathrm{unit}}_{\Lambda_{k+1}}
\in
[\Xi^{(1)}_{\Lambda_{k+1}}]\cap E^-_{\Lambda_{k+1}}.
\]
Then the following are equivalent:
\begin{enumerate}[leftmargin=2em,label=(\alph*)]
\item
  \[
  J^{\mathrm{unit}}_{\Lambda_{k+1}}
  -
  F_{\tau^{\mathrm{unit}}_{\Lambda_{k+1}},\Lambda_{k+1}}
  \in
  \mathcal A^-_{\Lambda_{k+1}};
  \]
\item
  \[
  q_{\Lambda_{k+1}}
  \!\bigl(
  J^{\mathrm{unit}}_{\Lambda_{k+1}}
  \bigr)
  =
  \overline\ell^{\mathrm{split,src}}_{\Lambda_{k+1}}.
  \]
\end{enumerate}
If these equivalent statements hold, then
\[
\overline{\rho}^{\mathrm{dom}}_{\Lambda_{k+1}}
\!\bigl(
q_{\Lambda_{k+1}}(J^{\mathrm{unit}}_{\Lambda_{k+1}})
\bigr)
=
c_{\Lambda_{k+1}}
\qquad\text{and}\qquad
\mathsf W^{\mathrm{red}}_{\Lambda_{k+1}}
\!\bigl(
J^{\mathrm{unit}}_{\Lambda_{k+1}}
\bigr)
=
c_{\Lambda_{k+1}}\,
P^{\mathrm{LoG}}_{-,\Lambda_{k+1}}.
\]
\end{proposition}

\begin{proof}
Statement \textnormal{(a)} is equivalent to
\[
q_{\Lambda_{k+1}}
\!\bigl(
J^{\mathrm{unit}}_{\Lambda_{k+1}}
\bigr)
=
q_{\Lambda_{k+1}}
\!\bigl(
F_{\tau^{\mathrm{unit}}_{\Lambda_{k+1}},\Lambda_{k+1}}
\bigr),
\]
which is exactly statement \textnormal{(b)} by the definition of
\(
\overline\ell^{\mathrm{split,src}}_{\Lambda_{k+1}}
\).
So the two formulations are equivalent.

If they hold, then
\[
q_{\Lambda_{k+1}}
\!\bigl(
J^{\mathrm{unit}}_{\Lambda_{k+1}}
\bigr)
=
\overline\ell^{\mathrm{split,src}}_{\Lambda_{k+1}}.
\]
Applying item \textnormal{(3)} of \cref{thm:splitlift-canonical-quotient-bridge} gives
\[
\overline{\rho}^{\mathrm{dom}}_{\Lambda_{k+1}}
\!\bigl(
q_{\Lambda_{k+1}}(J^{\mathrm{unit}}_{\Lambda_{k+1}})
\bigr)
=
c_{\Lambda_{k+1}}.
\]
Because
\(
J^{\mathrm{unit}}_{\Lambda_{k+1}}
\in
[\Xi^{(1)}_{\Lambda_{k+1}}]
\),
\cref{stack-cor:dominant-quotient-constancy} identifies the left-hand side with
\(
\rho^{\mathrm{dom}}_{\Lambda_{k+1}}(J^{\mathrm{unit}}_{\Lambda_{k+1}})
\).
Finally,
\cref{stack-prop:reduced-welded-ambiguity-trivial}
gives the displayed reduced welded image identity.
\end{proof}

\begin{theorem}[Canonical unit realization from one subordinate realization]
\label{thm:splitlift-canonical-unit-realization}
Fix an accepted tick \(k\) and the next accepted receiver step \(k+1\).
Assume:
\begin{enumerate}[leftmargin=2em]
\item the hypotheses of \cref{prop:splitlift-ambiguity-subordination};
\item there exists a source-compatible unit descendant realization
  \[
  J^{\mathrm{unit},0}_{\Lambda_{k+1}}
  \in
  [\Xi^{(1)}_{\Lambda_{k+1}}]\cap E^-_{\Lambda_{k+1}}
  \]
  satisfying the equivalent conditions of \cref{prop:splitlift-ambiguity-subordination};
\item the imported finite-cutoff canonical representative package at \(\Lambda_{k+1}\): there exists
  a Hilbertizing bilinear form
  \(
  \mathsf B_{\Lambda_{k+1}}
  \)
  on
  \(
  E^-_{\Lambda_{k+1}}
  \)
  and a unique representative
  \[
  \mathcal J^{\mathrm{can}}_{\Lambda_{k+1}}
  \in
  [\Xi^{(1)}_{\Lambda_{k+1}}]\cap E^-_{\Lambda_{k+1}}
  \]
  minimizing
  \(
  \mathsf B_{\Lambda_{k+1}}(J,J)
  \),
  equivalently satisfying
  \[
  \mathsf B_{\Lambda_{k+1}}
  \bigl(
  \mathcal J^{\mathrm{can}}_{\Lambda_{k+1}},
  A
  \bigr)
  =
  0
  \qquad
  \forall\,A\in\mathcal A^-_{\Lambda_{k+1}}.
  \]
\end{enumerate}
Set
\[
J^{\mathrm{unit,can}}_{\Lambda_{k+1}}
:=
\mathcal J^{\mathrm{can}}_{\Lambda_{k+1}}.
\]
Then:
\begin{enumerate}[leftmargin=2em]
\item
  \(
  J^{\mathrm{unit,can}}_{\Lambda_{k+1}}
  \in
  [\Xi^{(1)}_{\Lambda_{k+1}}]\cap E^-_{\Lambda_{k+1}}
  \);
\item
  \[
  q_{\Lambda_{k+1}}
  \!\bigl(
  J^{\mathrm{unit,can}}_{\Lambda_{k+1}}
  \bigr)
  =
  \overline\ell^{\mathrm{split,src}}_{\Lambda_{k+1}};
  \]
\item
  \[
  J^{\mathrm{unit,can}}_{\Lambda_{k+1}}
  -
  F_{\tau^{\mathrm{unit}}_{\Lambda_{k+1}},\Lambda_{k+1}}
  \in
  \mathcal A^-_{\Lambda_{k+1}};
  \]
\item
  \(
  J^{\mathrm{unit,can}}_{\Lambda_{k+1}}
  \)
  is canonical relative to
  \(
  \mathsf B_{\Lambda_{k+1}}
  \)
  and the source-side bridge point
  \(
  \overline\ell^{\mathrm{split,src}}_{\Lambda_{k+1}}
  \).
\end{enumerate}
Hence one subordinate unit realization is enough to promote the fixed quotient bridge point to a
canonical source-compatible unit descendant realization.
\end{theorem}

\begin{proof}
Item \textnormal{(1)} is immediate from the definition of
\(
J^{\mathrm{unit,can}}_{\Lambda_{k+1}}
\).
Because
\(
J^{\mathrm{unit,can}}_{\Lambda_{k+1}}
\)
and
\(
J^{\mathrm{unit},0}_{\Lambda_{k+1}}
\)
both lie in the same affine class
\(
[\Xi^{(1)}_{\Lambda_{k+1}}]
\),
their difference belongs to
\(
\mathcal A^-_{\Lambda_{k+1}}
\).
Since
\(
J^{\mathrm{unit},0}_{\Lambda_{k+1}}
\)
satisfies the equivalent conditions of \cref{prop:splitlift-ambiguity-subordination}, we have
\[
q_{\Lambda_{k+1}}
\!\bigl(
J^{\mathrm{unit},0}_{\Lambda_{k+1}}
\bigr)
=
\overline\ell^{\mathrm{split,src}}_{\Lambda_{k+1}}
\]
and
\[
J^{\mathrm{unit},0}_{\Lambda_{k+1}}
-
F_{\tau^{\mathrm{unit}}_{\Lambda_{k+1}},\Lambda_{k+1}}
\in
\mathcal A^-_{\Lambda_{k+1}}.
\]
Adding the ambiguity difference
\(
J^{\mathrm{unit,can}}_{\Lambda_{k+1}}
-J^{\mathrm{unit},0}_{\Lambda_{k+1}}
\in
\mathcal A^-_{\Lambda_{k+1}}
\)
to the second display proves item \textnormal{(3)}.
Applying the quotient map to the same ambiguity difference proves item \textnormal{(2)}.
Finally, item \textnormal{(4)} is exactly the uniqueness clause in the imported finite-cutoff
canonical-representative package.
\end{proof}

\begin{theorem}[Bridge closure from one ambiguity-subordinated unit realization]
\label{thm:splitlift-canonical-bridge-from-B}
Fix an accepted tick \(k\) and the next accepted receiver step \(k+1\).
Assume the hypotheses of \cref{thm:splitlift-canonical-unit-realization}.
Then the source-side quotient bridge point
\(
\overline\ell^{\mathrm{split,src}}_{\Lambda_{k+1}}
\)
is exactly the descendant-class quotient point
\(
q_{\Lambda_{k+1}}(J^{\mathrm{unit,can}}_{\Lambda_{k+1}})
\),
and the canonical map
\(
\mathfrak L^{\mathrm{split,src}}_{k\to k+1}
\)
is the desired split-to-descendant bridge datum in
\(
\mathfrak D^-_{\Lambda_{k+1}}
\).
\end{theorem}

\begin{proof}
This is exactly item \textnormal{(2)} of
\cref{thm:splitlift-canonical-unit-realization}, together with the definition of
\(
\mathfrak L^{\mathrm{split,src}}_{k\to k+1}
\)
from \cref{thm:splitlift-canonical-quotient-bridge}.
\end{proof}

\begin{corollary}[Subordinate unit realizations are unique modulo ambiguity]
\label{cor:splitlift-unit-realization-uniqueness}
Fix an accepted tick \(k\) and the next accepted receiver step \(k+1\).
Assume the hypotheses of \cref{prop:splitlift-ambiguity-subordination}.
Let
\[
J^{\mathrm{unit},1}_{\Lambda_{k+1}},
\quad
J^{\mathrm{unit},2}_{\Lambda_{k+1}}
\in
[\Xi^{(1)}_{\Lambda_{k+1}}]\cap E^-_{\Lambda_{k+1}}
\]
each satisfy the equivalent conditions of \cref{prop:splitlift-ambiguity-subordination}.
Then:
\begin{enumerate}[leftmargin=2em]
\item
  \[
  J^{\mathrm{unit},2}_{\Lambda_{k+1}}
  -
  J^{\mathrm{unit},1}_{\Lambda_{k+1}}
  \in
  \mathcal A^-_{\Lambda_{k+1}};
  \]
\item
  \[
  q_{\Lambda_{k+1}}
  \!\bigl(
  J^{\mathrm{unit},1}_{\Lambda_{k+1}}
  \bigr)
  =
  q_{\Lambda_{k+1}}
  \!\bigl(
  J^{\mathrm{unit},2}_{\Lambda_{k+1}}
  \bigr)
  =
  \overline\ell^{\mathrm{split,src}}_{\Lambda_{k+1}};
  \]
\item
  \[
  \overline{\rho}^{\mathrm{dom}}_{\Lambda_{k+1}}
  \!\bigl(
  q_{\Lambda_{k+1}}(J^{\mathrm{unit},1}_{\Lambda_{k+1}})
  \bigr)
  =
  \overline{\rho}^{\mathrm{dom}}_{\Lambda_{k+1}}
  \!\bigl(
  q_{\Lambda_{k+1}}(J^{\mathrm{unit},2}_{\Lambda_{k+1}})
  \bigr)
  =
  c_{\Lambda_{k+1}};
  \]
\item
  \[
  \mathsf W^{\mathrm{red}}_{\Lambda_{k+1}}
  \!\bigl(
  J^{\mathrm{unit},1}_{\Lambda_{k+1}}
  \bigr)
  =
  \mathsf W^{\mathrm{red}}_{\Lambda_{k+1}}
  \!\bigl(
  J^{\mathrm{unit},2}_{\Lambda_{k+1}}
  \bigr)
  =
  c_{\Lambda_{k+1}}\,
  P^{\mathrm{LoG}}_{-,\Lambda_{k+1}}.
  \]
\end{enumerate}
In particular, source-compatible unit realizations subordinate to the canonical source-side bridge
point are unique modulo
\(
\mathcal A^-_{\Lambda_{k+1}}
\),
and every datum read only through that quotient point together with the reduced welded image is
representative-independent on that class.
\end{corollary}

\begin{proof}
Because both
\(
J^{\mathrm{unit},1}_{\Lambda_{k+1}}
\)
and
\(
J^{\mathrm{unit},2}_{\Lambda_{k+1}}
\)
belong to
\(
[\Xi^{(1)}_{\Lambda_{k+1}}]
\),
their difference lies in
\(
\mathcal A^-_{\Lambda_{k+1}}
\),
which proves item \textnormal{(1)}.
Item \textnormal{(2)} is exactly the quotient-point clause of
\cref{prop:splitlift-ambiguity-subordination} applied to each representative.
Items \textnormal{(3)} and \textnormal{(4)} are then the final conclusions of
\cref{prop:splitlift-ambiguity-subordination}, again applied to each representative.
The concluding sentence is immediate from those displayed identities.
\end{proof}

\begin{remark}[Exact-family accepted-boundary invariance]
\label{rem:splitlift-unit-realization-accepted-boundary-audit}
The theorem-backed corollary above controls the quotient point and the reduced welded image.
On the current exact split-controlled surrogate family, the separate audit packet
\texttt{Main/Mathematica/Reports/Lane1BridgeClosureAudit.md} also certifies a stronger
finite-dimensional statement: any two admissible positive-calibration realizations subordinate to the
same source-side bridge point have the same accepted welded boundary image.
That accepted-boundary invariance is therefore presently exact-family support rather than
theorem-backed text.
\end{remark}

\begin{corollary}[Canonical representative reconstruction after bridge closure]
\label{cor:splitlift-canonical-bridge-reconstruction}
Fix an accepted tick \(k\) and the next accepted receiver step \(k+1\).
Assume the hypotheses of \cref{thm:splitlift-canonical-bridge-from-B}.
Set
\[
J^{\mathrm{split,can}}_{\Lambda_{k+1}}
:=
J^{\mathrm{unit,can}}_{\Lambda_{k+1}}
\]
and define
\[
\mathfrak I^{\mathrm{split}}_{\Lambda_{k+1}}
:
\mathfrak S^{\mathrm{split,can}}_{\Lambda_{k+1}}
\longrightarrow
E^-_{\Lambda_{k+1}},
\qquad
\mathfrak I^{\mathrm{split}}_{\Lambda_{k+1}}
\bigl(
a\,s^{\mathrm{split,can}}_{\Lambda_{k+1}}
\bigr)
:=
a\,J^{\mathrm{split,can}}_{\Lambda_{k+1}}.
\]
Then:
\begin{enumerate}[leftmargin=2em]
\item
  \[
  q_{\Lambda_{k+1}}
  \!\bigl(
  J^{\mathrm{split,can}}_{\Lambda_{k+1}}
  \bigr)
  =
  \overline\ell^{\mathrm{split,src}}_{\Lambda_{k+1}};
  \]
\item
  \[
  \mathsf W^{\mathrm{red}}_{\Lambda_{k+1}}
  \!\bigl(
  J^{\mathrm{split,can}}_{\Lambda_{k+1}}
  \bigr)
  =
  c_{\Lambda_{k+1}}\,
  P^{\mathrm{LoG}}_{-,\Lambda_{k+1}};
  \]
\item the pair
  \(
  \mathfrak T^{\mathrm{split}}_{k\to k+1}
  \)
  and
  \(
  \mathfrak I^{\mathrm{split}}_{\Lambda_{k+1}}
  \)
  defines a canonical split-to-descendant scalar bridge in the sense of
  \cref{def:splitlift-canonical-bridge}.
\end{enumerate}
\end{corollary}

\begin{proof}
Item \textnormal{(1)} is exactly \cref{thm:splitlift-canonical-bridge-from-B}.
Item \textnormal{(2)} is the reduced welded-image conclusion of
\cref{prop:splitlift-ambiguity-subordination}, applied to the canonical unit realization supplied by
\cref{thm:splitlift-canonical-unit-realization}.

By \cref{thm:splitlift-canonical-scalar-carrier}, the generator
\(
s^{\mathrm{split,can}}_{\Lambda_{k+1}}
\)
is exactly
\(
\widehat\ell^{\mathrm{delay}}_k
\)
and
\(
\mathfrak T^{\mathrm{split}}_{k\to k+1}
\)
is the identity on the delayed selected line.
So every clause of \cref{def:splitlift-canonical-bridge} is satisfied, proving item
\textnormal{(3)}.
\end{proof}

\begin{corollary}[\(B\)-localized recoherence plus one subordinate unit realization produces canonical scalar datum]
\label{cor:splitlift-canonical-datum-from-B}
Fix an accepted tick \(k\) and the next accepted receiver step \(k+1\).
Assume the hypotheses of \cref{cor:splitlift-canonical-bridge-reconstruction}.
Then \cref{prop:splitlift-canonical-datum} applies with the bridge reconstructed in
\cref{cor:splitlift-canonical-bridge-reconstruction}.
Consequently the split-controlled delayed selected-line data determine a canonical scalar datum
and a canonical quotient subspace
\[
\overline L^{\mathrm{Main,can}}_{\Lambda_{k+1}}
\subset
\mathfrak D^-_{\Lambda_{k+1}}.
\]
\end{corollary}

\begin{proof}
Apply \cref{prop:splitlift-canonical-datum} to the canonical bridge furnished by
\cref{cor:splitlift-canonical-bridge-reconstruction}.
\end{proof}

\begin{theorem}[Split-controlled scalar descendant lift]
\label{thm:splitlift-scalar-descendant-lift}
Fix an accepted tick \(k\) and the next accepted receiver step \(k+1\).
Assume:
\begin{enumerate}[leftmargin=2em]
\item the hypotheses and conclusions of \cref{thm:splitlift-main};
\item the delayed selected-line package of
  \cref{delay-def:delayed-selected-line,delay-def:delayed-response-map,delay-thm:delayed-response-sufficiency};
\item a split-controlled scalar descendant lift datum of
  \cref{def:splitlift-scalar-datum}.
\end{enumerate}
Set
\[
\overline\ell^{\mathrm{Main}}_{\Lambda_{k+1}}
:=
\mathfrak L^{\mathrm{split}}_{k\to k+1}
\bigl(\widehat\ell^{\mathrm{delay}}_k\bigr)
=
q_{\Lambda_{k+1}}\!\bigl(J^{\mathrm{split}}_{\Lambda_{k+1}}\bigr)
\]
and
\[
\overline L^{\mathrm{Main}}_{\Lambda_{k+1}}
:=
\mathbb R\,\overline\ell^{\mathrm{Main}}_{\Lambda_{k+1}}
\subset
\mathfrak D^-_{\Lambda_{k+1}}.
\]
Then:
\begin{enumerate}[leftmargin=2em]
\item the map
  \[
  \mathfrak L^{\mathrm{split}}_{k\to k+1}:
  \widehat L^{\mathrm{sel,delay}}_k\to \mathfrak D^-_{\Lambda_{k+1}}
  \]
  is well defined on the delayed selected line;
\item the image generator is nonzero:
  \[
  \overline\ell^{\mathrm{Main}}_{\Lambda_{k+1}}\neq 0;
  \]
\item the image generator is dominant-active:
  \[
  \overline{\rho}^{\mathrm{dom}}_{\Lambda_{k+1}}
  \bigl(
  \overline\ell^{\mathrm{Main}}_{\Lambda_{k+1}}
  \bigr)
  =
  c_{\Lambda_{k+1}}
  \neq 0;
  \]
\item \(\mathfrak L^{\mathrm{split}}_{k\to k+1}\) is faithful on
  \(\widehat L^{\mathrm{sel,delay}}_k\);
\item the current Main-visible reduced welded image of the lifted generator is nonzero:
  \[
  \mathsf W^{\mathrm{red}}_{\Lambda_{k+1}}
  \!\bigl(J^{\mathrm{split}}_{\Lambda_{k+1}}\bigr)
  \neq 0.
  \]
\end{enumerate}
\end{theorem}

\begin{proof}
By \cref{delay-thm:delayed-response-sufficiency}, the delayed realization map has trivial kernel on
\(
\widehat L^{\mathrm{sel,delay}}_k
\).
So the delayed selected line is already faithful on its own delayed carrier.

Because
\(
\widehat L^{\mathrm{sel,delay}}_k
=
\mathbb R\,\widehat\ell^{\mathrm{delay}}_k
\)
by \cref{delay-def:delayed-selected-line}, any linear map out of that line is determined uniquely by
the image of the generator.
The datum of \cref{def:splitlift-scalar-datum} therefore makes
\(
\mathfrak L^{\mathrm{split}}_{k\to k+1}
\)
well defined, proving item \textnormal{(1)}.

By the compatibility clause in \cref{def:splitlift-scalar-datum},
\[
\mathsf W^{\mathrm{red}}_{\Lambda_{k+1}}
\!\bigl(J^{\mathrm{split}}_{\Lambda_{k+1}}\bigr)
=
c_{\Lambda_{k+1}}\,
P^{\mathrm{LoG}}_{-,\Lambda_{k+1}}.
\]
Here
\(
c_{\Lambda_{k+1}}\neq 0
\)
by the delayed selected-line package and
\(
P^{\mathrm{LoG}}_{-,\Lambda_{k+1}}\neq 0
\)
by \cref{stack-cor:compressed-export-selection}.
Hence the displayed reduced welded image is nonzero, proving item \textnormal{(5)}.

Now suppose for contradiction that
\[
\overline\ell^{\mathrm{Main}}_{\Lambda_{k+1}}
=
q_{\Lambda_{k+1}}\!\bigl(J^{\mathrm{split}}_{\Lambda_{k+1}}\bigr)
=
0.
\]
Then
\(
J^{\mathrm{split}}_{\Lambda_{k+1}}\in \mathcal A^-_{\Lambda_{k+1}}
\),
so \cref{stack-prop:reduced-welded-ambiguity-trivial} gives
\[
\mathsf W^{\mathrm{red}}_{\Lambda_{k+1}}
\!\bigl(J^{\mathrm{split}}_{\Lambda_{k+1}}\bigr)
=
0,
\]
contradicting item \textnormal{(5)} already proved.
Therefore
\(
\overline\ell^{\mathrm{Main}}_{\Lambda_{k+1}}\neq 0
\),
which is item \textnormal{(2)}.

By \cref{stack-prop:reduced-welded-ambiguity-trivial}, the reduced welded image is exactly the
dominant-response scalar times the canonical LoG projector:
\[
\mathsf W^{\mathrm{red}}_{\Lambda_{k+1}}
\!\bigl(J^{\mathrm{split}}_{\Lambda_{k+1}}\bigr)
=
\rho^{\mathrm{dom}}_{\Lambda_{k+1}}
\!\bigl(J^{\mathrm{split}}_{\Lambda_{k+1}}\bigr)\,
P^{\mathrm{LoG}}_{-,\Lambda_{k+1}}.
\]
Comparing with the compatibility clause from \cref{def:splitlift-scalar-datum} shows
\[
\rho^{\mathrm{dom}}_{\Lambda_{k+1}}
\!\bigl(J^{\mathrm{split}}_{\Lambda_{k+1}}\bigr)
=
c_{\Lambda_{k+1}}.
\]
Because
\(
J^{\mathrm{split}}_{\Lambda_{k+1}}
\in
[\Xi^{(1)}_{\Lambda_{k+1}}]
\),
\cref{stack-cor:dominant-quotient-constancy} yields
\[
\overline{\rho}^{\mathrm{dom}}_{\Lambda_{k+1}}
\bigl(
\overline\ell^{\mathrm{Main}}_{\Lambda_{k+1}}
\bigr)
=
\rho^{\mathrm{dom}}_{\Lambda_{k+1}}
\!\bigl(J^{\mathrm{split}}_{\Lambda_{k+1}}\bigr)
=
c_{\Lambda_{k+1}}
\neq 0,
\]
which proves item \textnormal{(3)}.

Finally, the domain
\(
\widehat L^{\mathrm{sel,delay}}_k
\)
is one-dimensional, and its generator has nonzero image
\(
\overline\ell^{\mathrm{Main}}_{\Lambda_{k+1}}\neq 0
\).
Therefore the kernel of
\(
\mathfrak L^{\mathrm{split}}_{k\to k+1}
\)
is trivial, proving item \textnormal{(4)}.
\end{proof}

\begin{corollary}[Canonical split-controlled scalar descendant lift]
\label{cor:splitlift-canonical-scalar-descendant-lift}
Fix an accepted tick \(k\) and the next accepted receiver step \(k+1\).
Assume:
\begin{enumerate}[leftmargin=2em]
\item the hypotheses of \cref{cor:splitlift-canonical-datum-from-B};
\item the delayed selected-line package of
  \cref{delay-def:delayed-selected-line,delay-def:delayed-response-map,delay-thm:delayed-response-sufficiency}.
\end{enumerate}
Then the canonical datum of \cref{cor:splitlift-canonical-datum-from-B} satisfies the hypotheses of
\cref{thm:splitlift-scalar-descendant-lift}.
Consequently the conclusions of that theorem hold with
\[
J^{\mathrm{split}}_{\Lambda_{k+1}}
:=
J^{\mathrm{split,can}}_{\Lambda_{k+1}},
\qquad
\overline\ell^{\mathrm{Main}}_{\Lambda_{k+1}}
:=
\overline\ell^{\mathrm{Main,can}}_{\Lambda_{k+1}},
\]
and
\[
\overline L^{\mathrm{Main}}_{\Lambda_{k+1}}
=
\overline L^{\mathrm{Main,can}}_{\Lambda_{k+1}}
\subset
\mathfrak D^-_{\Lambda_{k+1}}.
\]
In particular, once the canonical split-to-descendant scalar bridge is supplied, the scalar datum of
the split-controlled delayed selected-line package is canonical and determines a canonical nonzero
dominant-active scalar line in the ambiguity quotient once the delayed selected-line package is in
force.
\end{corollary}

\begin{proof}
By \cref{cor:splitlift-canonical-datum-from-B}, the canonical bridge built from the
\(B\)-localized recoherence output induces a split-controlled scalar descendant lift datum.
The remaining assumptions of \cref{thm:splitlift-scalar-descendant-lift} are already included in
\cref{cor:splitlift-canonical-datum-from-B}, except for the delayed selected-line package, which is
assumed here.
So the conclusions of \cref{thm:splitlift-scalar-descendant-lift} apply verbatim to the canonical
datum.
\end{proof}

\begin{remark}[Boundary for canonical scalar data]
\label{rem:splitlift-canonical-status}
Theorem~\cref{thm:splitlift-canonical-quotient-bridge} determines the canonical source-side bridge
point
\(
\overline\ell^{\mathrm{split,src}}_{\Lambda_{k+1}}
\in
\mathfrak D^-_{\Lambda_{k+1}}
\)
before any descendant representative is chosen.
If there exists a source-compatible unit descendant realization subordinate to that bridge point,
then
Theorem~\cref{thm:splitlift-canonical-unit-realization} and
Corollary~\cref{cor:splitlift-unit-realization-uniqueness}
canonicalize that realization and reconstruct the canonical split-to-descendant scalar bridge.

What remains assumed is the existence of one such subordinate source-compatible unit descendant
realization.
Nor does the present note derive the package nondegeneracy clause
\(
c_{\Lambda_{k+1}}\neq 0
\);
that clause is still imported from the upstream seam-overlap/transport lane.
\end{remark}

\begin{remark}[Scope of the scalar bridge]
\label{rem:splitlift-scalar-bridge-scope}
\cref{thm:splitlift-scalar-descendant-lift} reaches only the ambiguity quotient
\(
\mathfrak D^-_{\Lambda_{k+1}}
\)
and the current Main-visible reduced welded image.
It does \emph{not} claim any automatic further reentry into a stronger selected descendant quotient,
physical quotient, or mod-null package.
That later step remains conditional.
\end{remark}

\section{Structural Limits}
\label{sec:splitlift-boundaries}

\begin{remark}[Conditional character]
\label{rem:splitlift-conditional}
\label{rem:splitlift-upgrade}
\label{rem:splitlift-gate}
The proofs in this note begin once the controlled completion branch decomposition of
\cref{ass:splitlift-controlled} is granted.
They show that the inherited ledger projectors are QND, the abstract accepted-record payload on
\(A\) is preserved exactly, and any recoherence certified here is confined to branchwise
completion data \(B\).
They do not derive the direct-sum realization
\(
\Hacc\cong\bigoplus_i (\mathbb C e_i\otimes \mathcal H_i^{(B)})
\)
or its branchwise controlled realization from microscopic UOM dynamics.
\end{remark}

\begin{remark}[Residual source-internal factorization problem]
\label{rem:splitlift-blocker}
\label{rem:splitlift-blocker-physical}
To remove the explicit split-data premise from \cref{thm:splitlift-main}, one would need a
source-internal statement on the accepted welded branch asserting that the compiled welded control is
determined entirely by the accepted-record sectors of the inherited ledger, so that the controlling
branch labels and branchwise completion actions are fixed by accepted-record data with no additional
completion-side choice.
That factorization problem is separate from the ledger-side consequences proved here.
In particular, the present theorem should not be read as identifying the full completion datum with
the current accepted-record sector.
\end{remark}

\begin{remark}[Gaussian-independent scope]
\label{rem:splitlift-gaussian-fence}
Nothing in this note proves Gaussian dilation existence, symplectic uniqueness, or completion
Hamiltonian reconstruction.
Those realization questions are not used in the proofs of
\cref{lem:splitlift-qnd,prop:splitlift-payload,cor:splitlift-classical,thm:splitlift-main}.
\end{remark}

\end{document}
